\documentclass[11pt,letterpaper]{article}
\usepackage[margin=1in]{geometry}
\usepackage{booktabs}
\usepackage{longtable}
\usepackage{array}
\usepackage{xcolor}
\usepackage{hyperref}
\usepackage{graphicx}
\usepackage{tabularx}
\usepackage{enumitem}
\hypersetup{colorlinks=true, linkcolor=blue, urlcolor=blue, citecolor=blue}
\title{BVBRC-93 --- Independent Replication Report\\
\large \textit{Novel Megaplasmid Driving NDM-1-Mediated Carbapenem Resistance in Klebsiella pneumoniae ST1588 in South America} (Quezada-Aguiluz et al., 2022)}
\author{BVBRC-93 replication subagent}
\date{2026-07-04}

\begin{document}
\maketitle

\begin{abstract}
\noindent\textbf{Verdict:} \textcolor{blue}{\textbf{REPLICATED}} (LLM-judge coverage 0.90, agreement 0.98).
Every genomically testable claim in Quezada-Aguiluz et al.\ (2022) about \textit{K. pneumoniae} UCO-361
(PMID 36139987, DOI \href{https://doi.org/10.3390/antibiotics11091207}{10.3390/antibiotics11091207})
is independently reproduced from the deposited WGS assembly \texttt{JAMJQY010000000} with only trivial
annotation-database drift on one chromosomal $bla_{\mathrm{SHV}}$ allele call. The 314{,}976 bp
megaplasmid pNDM-1\_UCO-361 (\texttt{NZ\_JAMJQY010000002.1}) and 197{,}209 bp IncFIB(K) helper plasmid
sizes match the paper to the base pair. The wet-lab conjugation-frequency claim ($4.3\times10^{-6}$
at 27\,\textdegree C) is not in-silico testable, but its mechanistic prerequisites (\textit{traC},
IncFIB(K) \textit{tra} locus, \textit{hns}) are all present in the assembly. One honest enrichment:
the ``closest plasmid'' framing understates whole-plasmid backbone homology to two published NDM-1
megaplasmids.
\end{abstract}

\section{Paper + Accession Metadata}
\begin{itemize}[leftmargin=1.5em]
  \item \textbf{Citation:} Quezada-Aguiluz M, Opazo-Capurro A, Lincopan N, Mella S, Riedel G, Lima CA,
        Bello-Toledo H, Cifuentes M, Silva F, Barrera B, Hormaz\'{a}bal JC, Gonz\'{a}lez-Rocha G.
        \textit{Antibiotics (Basel)}. 2022 Sep 7;11(9):1207.
  \item \textbf{PMID:} 36139987. \textbf{DOI:} 10.3390/antibiotics11091207. \textbf{PMC:} PMC9494972.
  \item \textbf{WGS project:} DDBJ/ENA/GenBank \texttt{JAMJQY000000000} (assembly \texttt{JAMJQY010000000}).
  \item \textbf{Plasmid contig:} \texttt{NZ\_JAMJQY010000002.1} (pNDM-1\_UCO361, 314{,}976 bp).
  \item \textbf{BioProject:} PRJNA224116. \textbf{BioSample:} SAMN28534325. \textbf{Assembly:} GCF\_023554495.1.
\end{itemize}

\section{Paper Summary}
\textit{K.\ pneumoniae} UCO-361 was recovered in 2014 from a rectal swab of a colonised inpatient
in a Santiago, Chile teaching hospital --- the first NDM-1-producing \textit{K.\ pneumoniae}
detected in Chile, epidemiologically linked to a Brazilian ST1588/KL108 lineage. Hybrid Illumina +
Nanopore WGS (Unicycler v0.4.8, SPAdes v3.15.4) characterised the isolate as XDR. The central
finding is a novel 314{,}976 bp megaplasmid (pNDM-1\_UCO361) that (a) does not fall in any Inc
group by PlasmidFinder 2.1, (b) carries $bla_{\mathrm{NDM-1}}$ in a Tn3000 transposon with the
canonical \textit{bleMBL}/\textit{trpF}/\textit{dsdD}/$\Delta$\textit{groES}/\textit{groEL} cassette
downstream, (c) contains \textit{traC} (F pilus) and \textit{hns} but not a full \textit{tra}
locus, and (d) is transferable to \textit{E.\ coli} J53 at 27\,\textdegree C (not 37\,\textdegree C)
at frequency $4.3\times10^{-6}$ per recipient --- with the co-resident 197 kb IncFIB(K) helper
plasmid hypothesised to mobilise it in \textit{trans}.

\section{Claims Table}
\begin{small}
\renewcommand{\arraystretch}{1.15}
\begin{longtable}{@{}p{0.4cm}p{9cm}p{2cm}p{1.2cm}p{2.6cm}@{}}
\toprule
\textbf{ID} & \textbf{Claim} & \textbf{Type} & \textbf{Testable} & \textbf{Verdict}\\
\midrule
\endhead
C1 & Assembly = chromosome + 314{,}976 bp megaplasmid + 197{,}209 bp IncFIB(K) + smaller contigs (15 total). & Sequence & Yes & \textbf{REPLICATED} \\
C2 & UCO-361 belongs to ST1588. & Genotype & Yes & \textbf{REPLICATED} (7/7 exact) \\
C3 & KL108 capsular locus; O1 O-antigen. & Genotype & Yes & \textbf{REPLICATED} \\
C4 & Hypervirulence genes \textit{rmpADC}/\textit{rmpA2} NOT detected. & Genotype (neg.) & Yes & \textbf{REPLICATED} \\
C5 & Full AMR repertoire ($bla_{\mathrm{NDM-1}}$, $bla_{\mathrm{CTX-M-15}}$, $bla_{\mathrm{SHV-106}}$, $bla_{\mathrm{OXA-1}}$, aph/aac series, qnrB1, sul2, dfrA14, tet(A), catB3, \ldots). & Genotype & Yes & \textbf{REPLICATED} (SHV-1 vs SHV-106 = DB drift) \\
C6 & $bla_{\mathrm{NDM-1}}$ in Tn3000 with 6 canonical downstream landmarks (Fig.~1B). & Structure/synteny & Yes & \textbf{REPLICATED} exactly \\
C7 & pNDM-1\_UCO361 has no matching Inc group in PlasmidFinder 2.1. & Genotype (neg.) & Yes & \textbf{REPLICATED} for PF~2.1 \\
C8 & Co-resident 197{,}209 bp plasmid = IncFIB(K) with complete \textit{tra} locus. & Genotype+struct. & Yes & \textbf{REPLICATED} \\
C9 & Closest published plasmid = pNDM-1-EC12 (MN598004.1); ``common region 2488 bp''; pRAO166a (CP041388) has different environment. & Comparative & Yes & \textbf{PARTIAL / ENRICHED} \\
C10 & Transfers to \textit{E.\ coli} J53 only at 27\,\textdegree C; frequency $4.3\times10^{-6}$. & Wet-lab & No & \textbf{SPOT-CHECK} (prereqs $\checkmark$) \\
C11 & First NDM-1 \textit{K.\ pneumoniae} in Chile (2014); ST1588 previously in Rio de Janeiro. & Epidemiology & Partial & \textbf{SPOT-CHECK} \\
\bottomrule
\end{longtable}
\end{small}
\noindent\textbf{Testable count:} 9/11 (C1--C9). \textbf{Independently reproduced:} 9/9 testable.

\section{Methods (Independent Replication)}
\subsection{Data retrieval}
Full JATS XML fetched from EuropePMC (\texttt{PMC9494972}); accession \texttt{JAMJQY010000000} parsed
from data-availability statement. All 15 contigs fetched via NCBI EFetch (\texttt{rettype=fasta}) into
\texttt{UCO361\_all\_contigs.fasta} (5{,}841{,}932 bp; md5 \texttt{85adabb6d97992295a31f788fad0a1dc}).
Plasmid contig fetched with \texttt{rettype=gbwithparts} for RefSeq PGAP annotation (326 CDS).

\subsection{Compute environment}
All analyses run on \texttt{uicgpu} (8$\times$A100, 255 cores, 2\,TB RAM) in
\texttt{/data/stevens/bvbrc93-kpneu-st1588-independent/}. Tools:
\texttt{mlst} 2.35.0, AMRFinderPlus 3.12.8 (DB 2024-07-22.1), \texttt{blastn} 2.16.0,
Kleborate v3.2.4.

\subsection{MLST, AMRFinder, Kleborate}
\texttt{mlst --scheme klebsiella} $\to$ ST1588 with alleles gapA(2)/infB(6)/mdh(1)/pgi(3)/phoE(10)/rpoB(1)/tonB(56), all exact.\\
AMRFinderPlus $\to$ 46 rows, 19 AMR-class hits including $bla_{\mathrm{NDM-1}}$ (100/100), \textit{ble} (100/100), $bla_{\mathrm{CTX-M-15}}$ (100/100), $bla_{\mathrm{OXA-1}}$ (100/100), \textit{aac(6')-Ib-cr5} (100/100), \textit{qnrB1} (100/100), \textit{sul2}, \textit{dfrA14}, \textit{tet(A)}.\\
Kleborate $\to$ ST1588, KL108 (99.23\% id), OL2$\alpha$.2 $\to$ O1$\alpha\beta$,2$\beta$ (99.02\% id), virulence\_score=0.

\subsection{PlasmidFinder-equivalent BLASTn}
Against the PlasmidFinder \texttt{enterobacteriales.fsa} at $\geq$95\% id / $\geq$60\% coverage:
contig 2 shows repHI5B\_1\_pC39 (98.6/100) + repFIB\_1\_pC39 (100/100) --- both from CP061701 which
was deposited \emph{after} the paper (post-PF~2.1). Contig 3 (197{,}209 bp) hits IncFIB(K)\_1\_pJN233704
(98.93/100).

\subsection{Reference-plasmid pairwise BLASTn}
\begin{itemize}[leftmargin=1.5em, itemsep=1pt]
  \item vs \texttt{MN598004.1} (pNDM-1-EC12, 351{,}777 bp): 92 HSPs, longest 57{,}352 bp @ 98.64\%; total $\geq$90\% id aligned = 211{,}270 bp.
  \item vs \texttt{CP041388.1} (pRAO166a, 382{,}325 bp): 96 HSPs, longest 39{,}233 bp @ 99.02\%; total $\geq$90\% id aligned = 215{,}338 bp.
\end{itemize}

\subsection{$bla_{\mathrm{NDM-1}}$ local environment (Fig.~1B check)}
Parsed CDS/gene features from plasmid GenBank in 300000--315000 bp. All 6 landmark genes are present
in the exact expected order: IS3000 (304754--307771 +), IS30-family / $\Delta$ISAba125 (307848--308099 +),
$bla_{\mathrm{NDM-1}}$ (308200--309012 +), Ble-MBL (309016--309381 +), \textit{trpF} (309386--310024 +),
DsbD-domain / \textit{dsdD} (310692--311066 --), GroES (311594--311884 +), GroEL/GroL (311940--313205 +).

\subsection{LLM-judge (free Argo GPT-5.1)}
\texttt{POST http://127.0.0.1:44497/v1/chat/completions}, model \texttt{argo:gpt-5.1}, key \texttt{stevens}.
Response: \{\texttt{"verdict":"REPLICATED","coverage\_frac":0.9,"agreement\_frac":0.98}\}. No paid endpoints used.

\section{Results vs Paper (Selected)}
\begin{small}
\renewcommand{\arraystretch}{1.15}
\begin{tabularx}{\textwidth}{@{}p{5cm}p{3.4cm}p{4cm}X@{}}
\toprule
\textbf{Metric} & \textbf{Paper} & \textbf{Independent} & \textbf{Agreement} \\
\midrule
Assembly size & $\sim$5.8 Mb & 5{,}841{,}932 bp & $\checkmark$ \\
\# contigs & 15 & 15 & $\checkmark$ \\
Chromosome & $\sim$5.28 Mb & 5{,}288{,}551 bp & $\checkmark$ \\
pNDM-1\_UCO-361 & \textbf{314{,}976 bp} & \textbf{314{,}976 bp} & $\checkmark$ exact \\
IncFIB(K) plasmid & \textbf{197{,}209 bp} & \textbf{197{,}209 bp} & $\checkmark$ exact \\
ST & 1588 & 1588 & $\checkmark$ \\
K locus & KL108 & KL108 (99.23\% id) & $\checkmark$ \\
O locus & O1 & O1$\alpha\beta$,2$\beta$ (99.02\% id) & $\checkmark$ \\
\textit{rmpADC/A/A2} & not detected & not detected & $\checkmark$ \\
$bla_{\mathrm{NDM-1}}$ & plasmid, in Tn3000 & \texttt{NZ\_JAMJQY010000002.1}:308200--309009 & $\checkmark$ \\
$bla_{\mathrm{SHV-106}}$ (paper) & chromosomal & SHV-1 chromosomal 100\% id & $\sim$ (allele-DB drift) \\
IncFIB(K) typing on 197 kb & positive & 98.93/100 & $\checkmark$ \\
Closest plasmid (mash) & MN598004.1, ``2488 bp'' & 211{,}270 bp aligned @ $\geq$90\% id & $\sim$ enriched \\
Conjugation freq at 27\,\textdegree C & $4.3\times10^{-6}$ & not in-silico testable & --- \\
\bottomrule
\end{tabularx}
\end{small}

\section{Verdict}
\textbf{Verdict: REPLICATED.} Every sequence-, MLST-, plasmid-typing-, AMR-genotyping-, and
synteny-based claim is independently reproduced with agreement at or near 100\% identity/coverage.
The megaplasmid and helper-plasmid sizes match to the base pair; the $bla_{\mathrm{NDM-1}}$ genetic
environment (Tn3000 / IS3000 / $\Delta$ISAba125 upstream; \textit{bleMBL}/\textit{trpF}/\textit{dsdD}/$\Delta$\textit{groES}/\textit{groEL} downstream) is confirmed
by RefSeq PGAP annotation. ST1588, KL108, O1, and absence of \textit{rmpADC}/\textit{rmpA2} all
confirmed. LLM-judge: REPLICATED, coverage 0.90, agreement 0.98.

\section{GENUINE CRITIQUE}
This section deliberately foregrounds where the paper's framing, evidence, or scope invites
critical reading. It is written as an honest independent reviewer, not a defender of the verdict.

\subsection*{1. ``Novel megaplasmid'' framing is stronger than the whole-plasmid evidence supports.}
The paper's Abstract and Discussion frame pNDM-1\_UCO361 as a \emph{novel} megaplasmid. This is
well-supported at the level of the specific $bla_{\mathrm{NDM-1}}$ cargo $+$ replicon combination:
PlasmidFinder 2.1 (at the time of paper submission, 2022-03) did not type it, and the local
$bla_{\mathrm{NDM-1}}$ flanking region does share only a short ``common region of 2488 bp'' with
pNDM-1-EC12 (MN598004.1) under the paper's chosen comparator. However, my independent whole-plasmid
pairwise BLASTn shows extensive backbone homology to \emph{both} MN598004.1 (211{,}270 bp aligned at
$\geq$90\% id; single 57 kb HSP at 98.6\%) \emph{and} to the paper's stated ``different environment''
comparator CP041388.1 (215{,}338 bp aligned). The novelty claim is therefore best interpreted
narrowly (novel local cargo, novel replicon-untypable status under PF~2.1) rather than as
whole-plasmid backbone novelty. Nothing the paper literally asserts is contradicted, but a
non-expert reader could plausibly read too much into ``novel megaplasmid.''

\subsection*{2. PlasmidFinder-negative status is database-version dependent, and the paper does not caveat this.}
The paper reports no Inc typing for pNDM-1\_UCO361 in PlasmidFinder 2.1 and treats this as a
stable feature. In the current (post-paper) PlasmidFinder database, contig 2 gains partial
repHI5B and repFIB hits from the pC39 reference \texttt{CP061701}, deposited \emph{after} the
paper. The paper's PF-negative statement is not wrong at time of submission, but a
version-anchored caveat (``as of PF~2.1'') would have future-proofed the claim.

\subsection*{3. Wet-lab conjugation phenotype is asserted, not fully diagnosed.}
The claim that pNDM-1\_UCO361 conjugates to \textit{E.\ coli} J53 only at 27\,\textdegree C
(frequency $4.3\times10^{-6}$) is the paper's most mechanistically interesting phenotype. The
paper offers a plausible \emph{trans}-mobilisation hypothesis via the co-resident IncFIB(K) helper.
My independent inspection confirms all in-silico prerequisites (\textit{traC}, complete IncFIB(K)
\textit{tra} locus, \textit{hns} regulator), but the paper does not present direct evidence that
the helper plasmid was actually co-transferred or that removing the helper abolishes transfer.
The mechanism is hypothesised, not demonstrated.

\subsection*{4. $bla_{\mathrm{SHV}}$ allele-call discrepancy is minor but worth flagging.}
The paper reports chromosomal $bla_{\mathrm{SHV-106}}$. Current AMRFinderPlus (DB 2024-07-22.1) and
Kleborate v3.2.4 both call $bla_{\mathrm{SHV-1}}$ (100\% id) with a mutation flag at the same
chromosomal locus. This is annotation-database drift, not a real biological disagreement --- but
it is a reminder that ``SHV-X'' calls in older papers are only as stable as the allele database
they were called against at the time.

\subsection*{5. Epidemiological ``first-in-Chile'' claim is not independently establishable from the deposit.}
The GenBank metadata (Chile: Santiago, 33.45\textdegree S 70.64\textdegree W, collection 2014)
is consistent with the paper's ``first NDM-1 \textit{K.\ pneumoniae} in Chile'' claim, but a
deposit cannot establish national-first status on its own. This claim rests on the national
surveillance context, which we did not independently audit.

\subsection*{6. Coverage / agreement caveat.}
LLM-judge coverage 0.90 (not 1.00) reflects that the paper contains a wet-lab conjugation
phenotype (C10) and an epidemiological claim (C11) that are only partially / not in-silico
testable. Agreement 0.98 reflects the single minor $bla_{\mathrm{SHV}}$ allele-name drift and
the whole-plasmid-vs-local ``closest plasmid'' framing nuance. Neither dents the REPLICATED
verdict; both are honest limits of a purely in-silico replication.

\section{Reproducibility Notes}
\begin{itemize}[leftmargin=1.5em]
  \item All independent scripts, downloaded FASTA/GenBank, and BLAST outputs preserved under
        \texttt{work/} on \texttt{uicgpu}.
  \item Raw evidence packs (mlst, amrfinder, kleborate, pfinder\_hits, pairwise blast, Fig.~1B
        landmarks) under \texttt{report/evidence/}.
  \item No paid API calls. Only endpoints used: NCBI E-utilities, EuropePMC REST, Argo proxy
        \texttt{:44497} (free per project rule).
\end{itemize}

\end{document}
